% !PRINTED 302
\weberpage{302}

%<<CONT
es auch eine unter ihnen geben, die, ohne daß $r, s$ verschwinden, mit $V$ identisch wird. Dies verlangt aber

\begin{align*}
\alpha\equiv\pm\alpha g^{r+s}, &\qquad \beta\equiv\pm\beta g^{r-s}\\
\gamma\equiv\pm\gamma g^{-r+s}, &\qquad \delta\equiv\pm\delta g^{-r-s}\quad(\mathrm{mod}\ p),
\end{align*}

wo in allen vier Formeln die oberen oder die unteren Zeichen gelten, also, da weder $\alpha$ und $\delta$, noch $\beta$ und $\gamma$ gleichzeitig verschwinden,

\[ r+s\equiv 0,\qquad r-s\equiv 0\qquad(\mathrm{mod}\ p-1), \]
oder
\[ r+s\equiv\frac{p-1}{2},\qquad r-s\equiv\frac{p-1}{2}\qquad(\mathrm{mod}\ p-1), \]

also in beiden Fällen

\[ r\equiv 0,\qquad s\equiv 0\ \left(\mathrm{mod}\ \frac{p-1}{2}\right),\qquad\text{w. z. b. w.} \]

In der Form (16) sind also $\left(\dfrac{p-1}{2}\right)^2$ \gesperrt{verschiedene Substitutionen enthalten}. Ist damit die Gruppe $\mathfrak{K}$ noch nicht erschöpft, so wähle man eine nicht in (13), (14), (16) enthaltene Substitution $V'$ und bilde in gleicher Weise die Reihe

\weq{17}{C^{r}V'C^{s},}

deren Substitutionen sowohl unter sich als auch von den in (13), (14), (16) enthaltenen verschieden sind, und fahre auf diese Weise fort, bis die ganze Gruppe $\mathfrak{K}$ erschöpft ist.

Bezeichnen wir die Anzahl der so gebildeten Reihen (16), (17), $\dots$ mit $q$, so ergibt sich also der Grad der Gruppe $\mathfrak{K}$:

1. wenn $B$ in $\mathfrak{K}$ nicht vorkommt:

\[ =\frac{p-1}{2}+q\left(\frac{p-1}{2}\right)^2; \]

2. wenn $B$ in $\mathfrak{K}$ vorkommt:

\[ =p-1+q\left(\frac{p-1}{2}\right)^2, \]

und diese Zahl soll also nach der Forderung unserer Aufgabe

\[ =\frac{(p-1)(p+1)}{2} \]

sein. Hieraus aber ergibt sich, daß der Fall 1. unmöglich ist, denn es müßte in diesem Falle

\[ q\left(\frac{p-1}{2}\right)=p, \]
%CONT>>
